\scanpage{16}
% Reused typescript, rotated 180 degrees in the scan.
\figure{\begin{tikzpicture}\draw(0,0)--(.35,-.03);\draw(.65,-.06)--(.9,-.1);\end{tikzpicture}}
\begin{center}-2.29-\end{center}

\uline{Proof}: Immediate, since $\overline{\lambda_1\circ\lambda_2}=\bar\lambda_1+\bar\lambda_2$.

\uline{Corollary 2.16}: Let a reductive $G$ act on $X$, proper over $k$. Let $L\in\operatorname{Pic}^G(X)$ be ample, and let $x\in X_k$. Then $\{\delta\mid\nu^L(x,\delta)\leq0\}$ is semi-convex, and $\{\delta\mid\nu^L(x,\delta)<0\}$ is convex.

\uline{Proof}: Immediate.

It follows that from this Corollary and from Theorem 2.1, that if $\unsure{x}\in X_k$ is not \uline{semi-stable}, then we can associate to $x$ the convex set $C$ of $\delta\in\Delta(G)$ for which $\nu^L(x,\delta)<0$. Roughly speaking, we can say that a parabolic subgroup $P$ is ``responsible'' for this break-down of semi-stability if $C\cap\Delta_P(G)\neq\unsure{\phi|\varnothing}$. Since $C$ is convex, not too many parabolic subgroups $P$ are responsible; according to a conjecture of Tits, there is \inserted{even} a natural way to find one $P$ which is \uline{most} responsible, i.e.\ $P(\delta)$ for $\delta=$ the ``center'' of $C$ (cf.\ Preface).
Another consequence of Proposition 2.14 is:

\uline{Proposition 2.17}: Let a reductive $G$ act via $\sigma$ on $X$, proper over $k$. Let $L\in\operatorname{Pic}^G(X)$ be ample. Then there is a constant $K$ such that
\[
|\nu^L(x,\delta)|\leq K
\]
for all $x\in X_k$ and $\delta\in\Delta(G)$. Moreover, for every subset $S\unsure{\subset}X_k$, there is a $\delta_0\in\Delta(G)$ and $x_0\in S$ such that
\[
\nu^L(x_0,\delta_0)\leq\nu^L(x,\delta)^{\unsure{\circ|0}}
\]
for every other $\delta\in\Delta(G)$, $x\in S$.
\mspage{15}
